% =============================================================================
%  EE 362 Computer Organization Laboratory - report preamble
%  You do not need to edit this file. Everything here is set up for you.
% =============================================================================

\usepackage[utf8]{inputenc}
\usepackage[T1]{fontenc}
\usepackage{lmodern}
\usepackage[margin=2.4cm]{geometry}
\usepackage{graphicx}
\usepackage{booktabs}
\usepackage{amsmath}
\usepackage{amssymb}
\usepackage{enumitem}
\usepackage{xcolor}
\usepackage{listings}
\usepackage{float}
\usepackage{fancyhdr}
\usepackage{caption}
\usepackage[hidelinks]{hyperref}

% --- Colours -----------------------------------------------------------------
\definecolor{nmsucrimson}{RGB}{140,21,42}
\definecolor{codebg}{RGB}{248,248,248}
\definecolor{codekw}{RGB}{0,90,160}
\definecolor{codecomment}{RGB}{100,110,100}
\definecolor{codestring}{RGB}{160,60,20}

% --- VHDL code listings ------------------------------------------------------
\lstdefinestyle{vhdlstyle}{
  language=VHDL,
  backgroundcolor=\color{codebg},
  basicstyle=\ttfamily\scriptsize,
  keywordstyle=\color{codekw}\bfseries,
  commentstyle=\color{codecomment}\itshape,
  stringstyle=\color{codestring},
  numbers=left,
  numberstyle=\tiny\color{gray},
  numbersep=7pt,
  frame=single,
  framerule=0.4pt,
  rulecolor=\color{gray!40},
  breaklines=true,
  breakatwhitespace=true,
  showstringspaces=false,
  tabsize=4,
  columns=fullflexible,
  keepspaces=true,
  captionpos=b,
  xleftmargin=14pt,
}
\lstset{style=vhdlstyle}

% --- One-line macro for including a real .vhd file ---------------------------
%  Usage:  \vhdlfile{code/lab01/TransLatch.vhd}{TransLatch: 1-bit latch}
%
%  ALWAYS include code this way. Never paste code into the .tex file - pasted
%  code drifts out of date the moment you fix a bug, and graders can tell.
\newcommand{\vhdlfile}[2]{%
  \lstinputlisting[caption={#2},label={lst:\detokenize{#1}}]{#1}%
}

% --- Figure helper -----------------------------------------------------------
%  Usage: \resultfig{figs/lab01/waveform.png}{What this proves.}{fig:lab01-wave}
%
%  The second argument is NOT optional and is NOT a title. It is the sentence
%  that says what the figure proves. A figure without one earns zero marks.
\newcommand{\resultfig}[3]{%
  \begin{figure}[htbp]
    \centering
    \includegraphics[width=\linewidth]{#1}
    \caption{#2}
    \label{#3}
  \end{figure}%
}

% --- Placeholder, so the template compiles before you have any screenshots ---
%  DELETE these as you replace them with \resultfig. Any \figplaceholder left
%  in your final report is a missing figure.
\newcommand{\figplaceholder}[1]{%
  \begin{figure}[H]
    \centering
    \fbox{\parbox[c][3.2cm][c]{0.86\linewidth}{\centering\itshape
      PLACEHOLDER --- insert #1 here, then replace this command with
      \texttt{\textbackslash resultfig}}}
    \caption{Placeholder: #1.}
  \end{figure}%
}

% --- Headers -----------------------------------------------------------------
\pagestyle{fancy}
\fancyhf{}
\fancyhead[L]{\small EE\,362 Laboratory Report}
\fancyhead[R]{\small\leftmark}
\fancyfoot[C]{\small\thepage}
\renewcommand{\headrulewidth}{0.4pt}

\captionsetup{font=small,labelfont=bf}
\setlength{\parskip}{4pt}
